\FloatBarrier
\section{Limitations}
\label{sec:limitations}
The sensitivity range must be justified for the target population and covariate resolution. Statistical certificates additionally require overlap, controlled score complexity, and valid nuisance-error bounds. Small tail mass and scarce observed labels can make the bounds conservative. For \(\Gamma>1\), outcome-regression error contributes quadratically even with a correct observation model. Choosing \(\Gamma\) from data or examining multiple sensitivity levels requires corresponding joint error control (Appendix~\ref{app:gamma-calibration}).

The historical applicant analysis depends on regression extrapolation into unsupported regions and
on differently measured source fields. Controlled masking checks known issued-loan outcomes.
Repayment under an untaken lending action requires a causal or experimental model beyond passive
recovery of existing outcomes.

The revelation optimum fixes the predictor and threshold and assumes componentwise retained
sensitivity bounds. Known outcome-independent audit randomization can identify more, and direct
minimization of the reoptimized interval can favor a different allocation. Finite audits must
estimate their newly known contributions. Learned acquisition scores, retraining, and data-dependent
stopping require uncertainty control across the resulting audit history; the sequential trajectories
provide empirical performance evidence. Extensions to richer outcomes, unbounded losses, and other
application domains require additional modeling and validation.
